\subsection*{死锁}

默认情况下，事件是无锁的，用于不同异步块之间的高性能同步。这意味着多个线程同时更新事件时有风险，尤其是 \co@wait@。

例如：

\begin{lstlisting}[language=croqtile,caption={单实例死锁}]
parallel p by 1 {
  shared event e;
  trigger e;   // one thread set e;
  parallel q by 2 {
    wait e;    // two threads can reset e;
  }
}
\end{lstlisting}

这段代码里，两个线程等待同一个事件 \co@e@。但 \co@wait@ 会复位事件 \co@e@，因此可能出现 \co@thread-0@ 在 \co@thread-1@ 发现 \co@e@ 已触发之前就把它复位的情况，结果 \co@thread-1@ 永远无法解除阻塞，形成死锁。

\begin{lstlisting}[language=croqtile,caption={信用计数死锁}]
qc = 2;
parallel p by 1 {
  shared event e[qc];
  trigger e;   // one thread set e[0] and e[1];
  parallel q by qc {
    wait e[q]; // no dead lock
  }
}
\end{lstlisting}

这段代码使用事件数组，保证每个事件只由一个线程等待和复位。
